\section*{References}

\subsection*{Knowledge Base References}
\renewcommand{\refname}{}
\vspace{-1.2\baselineskip}
\begin{thebibliography}{KB99}

\bibitem[KB1]{fabryFeedbackHandbook}
G. Fabry.
\newblock \emph{Handbuch Feedback: Formatives Feedback in Studium und Lehre der Humanmedizin}.
\newblock Teaching handbook included in the Examiner Coach knowledge base, n.d.

\bibitem[KB2]{gigante2011getting}
J. Gigante, M. Dell, and A. Sharkey.
\newblock Getting beyond ``good job'': How to give effective feedback.
\newblock \emph{Pediatrics}, 127(2):205--210, 2011.
\newblock doi:10.1542/peds.2010-3351.

\bibitem[KB3]{hattie2007power}
J. Hattie and H. Timperley.
\newblock The power of feedback.
\newblock \emph{Review of Educational Research}, 77(1):81--112, 2007.
\newblock doi:10.3102/003465430298487.

\bibitem[KB4]{hewson1998giving}
M. G. Hewson and M. L. Little.
\newblock Giving feedback in medical education: Verification of recommended techniques.
\newblock \emph{Journal of General Internal Medicine}, 13(2):111--116, 1998.
\newblock doi:10.1046/j.1525-1497.1998.00027.x.

\bibitem[KB5]{hillienhof2018feedback}
A. Hillienhof.
\newblock Feedback statt Noten.
\newblock \emph{Medizin Studieren}, 2018.
\newblock Interview/article included in the Examiner Coach knowledge base.

\bibitem[KB6]{kunz2019frequency}
K. Kunz, M. Burkert, F. Heindl, K. Sch{\"u}ttpelz-Brauns, and M. Giesler.
\newblock The frequency of using certain feedback methods in the teaching of medicine: A survey of teachers at the medical faculties in Baden-Wurttemberg.
\newblock \emph{GMS Journal for Medical Education}, 36(4):Doc45, 2019.
\newblock doi:10.3205/zma001253.

\bibitem[KB7]{medic2022mentoringFlyer}
Medizincampus Chemnitz and Medizinische Fakult{\"a}t Carl Gustav Carus, Technische Universit{\"a}t Dresden.
\newblock \emph{Mentoring: Modellstudiengang Humanmedizin MEDiC}.
\newblock Flyer included in the Examiner Coach knowledge base, 2022.

\bibitem[KB8]{medic2023mentoringGuide}
Medizincampus Chemnitz and Medizinische Fakult{\"a}t Carl Gustav Carus, Technische Universit{\"a}t Dresden.
\newblock \emph{Leitfaden f{\"u}r das 1:1-Mentoring}.
\newblock Version dated September 2023.

\bibitem[KB9]{mitzFeedbackPoster}
MITZ.
\newblock \emph{Feedback}.
\newblock Poster included in the Examiner Coach knowledge base, 2018.

\bibitem[KB10]{ramani2012twelve}
S. Ramani and S. K. Krackov.
\newblock Twelve tips for giving feedback effectively in the clinical environment.
\newblock \emph{Medical Teacher}, 34(10):787--791, 2012.
\newblock doi:10.3109/0142159X.2012.684916.

\bibitem[KB11]{thrien2020feedback}
C. Thrien, G. Fabry, A. H{\"a}rtl, C. Kiessling, T. Graupe, I. Preusche, S. Pruskil, K. P. Schnabel, M. Sennekamp, S. R{\"u}ttermann, and A. W{\"u}nsch.
\newblock Feedback in medical education: A workshop report with practical examples and recommendations.
\newblock \emph{GMS Journal for Medical Education}, 37(5):Doc46, 2020.
\newblock doi:10.3205/zma001339.

\bibitem[KB12]{tuDresden2026feedbackBasics}
Medizinische Fakult{\"a}t Carl Gustav Carus, Technische Universit{\"a}t Dresden.
\newblock \emph{Feedback in der Lehre -- Basics: Allgemeines und Feedbackregeln}.
\newblock Internal teaching material included in the Examiner Coach knowledge base, 2026.

\bibitem[KB13]{vanDeRidder2023pendleton}
J. M. M. van de Ridder and M. Wijnen-Meijer.
\newblock Pendleton's rules: A mini review of a feedback method.
\newblock \emph{American Journal of Biomedical Science \& Research}, 19(1):19--21, 2023.
\newblock doi:10.34297/AJBSR.2023.19.002542.

\bibitem[KB14]{wisniewski2020power}
B. Wisniewski, K. Zierer, and J. Hattie.
\newblock The power of feedback revisited: A meta-analysis of educational feedback research.
\newblock \emph{Frontiers in Psychology}, 10:3087, 2020.
\newblock doi:10.3389/fpsyg.2019.03087.

\end{thebibliography}

\subsection*{Technical References}
\renewcommand{\refname}{}
\vspace{-1.2\baselineskip}
\begin{thebibliography}{T99}

\bibitem[T1]{gao2024modularrag}
Y. Gao, Y. Xiong, M. Wang, and H. Wang.
\newblock Modular RAG: Transforming RAG systems into LEGO-like reconfigurable frameworks.
\newblock arXiv:2407.21059, 2024.
\newblock \url{https://arxiv.org/abs/2407.21059}.

\bibitem[T2]{huang2025gumbel}
S. Huang, Z. Ma, J. Du, C. Meng, W. Wang, J. Leng, M. Guo, and Z. Lin.
\newblock Gumbel reranking: Differentiable end-to-end reranker optimization.
\newblock arXiv:2502.11116, 2025.
\newblock \url{https://arxiv.org/abs/2502.11116}.

\bibitem[T3]{jeong2024adaptiverag}
S. Jeong, J. Baek, S. Cho, S. J. Hwang, and J. C. Park.
\newblock Adaptive-RAG: Learning to adapt retrieval-augmented large language models through question complexity.
\newblock arXiv:2403.14403, 2024.
\newblock \url{https://arxiv.org/abs/2403.14403}.

\bibitem[T4]{jiang2024piperag}
W. Jiang, S. Zhang, B. Han, J. Wang, B. Wang, and T. Kraska.
\newblock PipeRAG: Fast retrieval-augmented generation via algorithm-system co-design.
\newblock arXiv:2403.05676, 2024.
\newblock \url{https://arxiv.org/abs/2403.05676}.

\bibitem[T5]{jin2025longcontext}
B. Jin, J. Yoon, J. Han, and S. O. Arik.
\newblock Long-context LLMs meet RAG: Overcoming challenges for long inputs in RAG.
\newblock In \emph{The Thirteenth International Conference on Learning Representations}, 2025.
\newblock \url{https://openreview.net/forum?id=oU3tpaR8fm}.

\bibitem[T6]{li2025enhancing}
S. Li, L. Stenzel, C. Eickhoff, and S. A. Bahrainian.
\newblock Enhancing retrieval-augmented generation: A study of best practices.
\newblock arXiv:2501.07391, 2025.
\newblock \url{https://arxiv.org/abs/2501.07391}.

\bibitem[T7]{wang2025diversity}
Z. Wang, B. Bi, Y. Luo, S. Asur, and C. N. Cheng.
\newblock Diversity enhances an LLM's performance in RAG and long-context task.
\newblock arXiv:2502.09017, 2025.
\newblock \url{https://arxiv.org/abs/2502.09017}.

\bibitem[T8]{kisski2026mission}
KISSKI.
\newblock KISSKI -- AI Service Centre for Sensitive and Critical Infrastructure.
\newblock 2026.
\newblock \url{https://kisski.gwdg.de/en/}. Accessed 8 June 2026.

\end{thebibliography}
